\section{The distinct squared-unit cubic quotient}
\label{z2x-section}

Retain \(\zeta^2-\zeta+1=0\) and the projective conventions of
Sections~\ref{rl-section} and~\ref{hg-section}. Put
\[
 A_0=S^3-3ST^2-T^3,\qquad B_0=3ST(S+T).
\]
The actual multiplication identity is
\[
 \zeta^2(S+T\zeta)^3=-(A_0+B_0)+A_0\zeta .
\]
Thus the coefficient \(\zeta^2\) gives \(r=-(A_0+B_0)/A_0\).
It is treated separately from the coefficient \(\zeta\); conjugation
does not preserve the required positive interval for \(r\).

\begin{proposition}[The first quotient and the common-source square gate]
\label{z2x-model}
The first hyperelliptic quotient of \(D_{3,\zeta^2}\) is the smooth
projective curve
\begin{equation}
 H:y^2=-3s(s+1)(s^3-3s-1),\qquad s=S/T.
 \label{z2x-quintic}
\end{equation}
It has genus two and the rational branch points \(P_0,P_{-1},P_\infty\)
over \(s=0,-1,\infty\). On its nonbranch chart let \(v=y/A_0(s,1)\).
The rational points that lift to the actual positive domain are
exactly those for which
\begin{equation}
 1<v<2/\sqrt3,\qquad 4-3v^2\text{ is a square in }\mathbb Q .
 \label{z2x-positive}
\end{equation}
Both signs of its nonzero square root \(w\) are allowed.
\end{proposition}
\begin{proof}
The equation \(v^2=1+r=-B_0/A_0\), after multiplication by \(A_0^2\),
gives the homogeneous equation \(Y^2=-A_0B_0\), where \(Y\) has
weight three. The cubic \(s^3-3s-1\) has discriminant \(81\) and no
rational root: neither of the possible integral roots \(1,-1\)
vanishes. Its values at \(0,-1\) are \(-1,1\). The quintic is
therefore squarefree, and its smooth projective model has genus
two and one branch point at infinity. The zeros of \(A_0\) are
nonrational cubic sources. They are retained projectively, although
none can be the source of a rational point.

Every rational point of \(D_{3,\zeta^2}\) maps to \(H\). At the same
source, the second square is exactly
\[
 w^2=1-3r=4+3B_0/A_0=4-3v^2.
\]
Condition~\eqref{z2x-positive} gives \(0<r=v^2-1<1/3\) and
\[
 t=\frac{v+w+2}{v-w}>1.
\]
Indeed \(v>1>|w|\), and the numerator exceeds the positive
denominator by \(2(w+1)>0\). Conversely, the positive domain in
Section~\ref{hg-section} gives precisely these conditions. Its
two signs of \(w\) give the ordered seed orientations. Both square
roots are nonzero, so this chart is smooth and normalization
introduces no further branch there.

At the three rational branch sources, \(B_0=0\), \(A_0\ne0\)
homogeneously, and \(v=0,r=-1\). They lift with \(w=\pm2\).
The original conic parameter is \(t=-2\) for \(w=2\), and \(t=0\)
for \(w=-2\), by direct substitution in the displayed inverse or
the original conic formulas. All lie outside
\eqref{z2x-positive}. No exhaustive rational-point claim is used
in this argument.
\end{proof}

\begin{proposition}[A rational cyclic cubic cover]
\label{z2x-cubic-map}
The rational function \(v=y/A_0\) defines a Galois morphism
\(H\to\mathbb P^1_v\) of degree three over \(\mathbb Q\).
Its generator is
\begin{equation}
 \rho(s,y)=\left(-\frac1{s+1},-\frac{y}{(s+1)^3}\right).
 \label{z2x-rho}
\end{equation}
There are four branch values, the simple zeros of \(v^4-v^2+1\),
each of ramification index three. Both \(v=0\) and \(v=\infty\)
are unramified.
\end{proposition}
\begin{proof}
The homogeneous substitution \((S,T)\mapsto(-T,S+T)\) sends
both \(A_0,B_0\) to their negatives. It preserves \(-A_0B_0\)
and induces~\eqref{z2x-rho}. The matrix has cube \(-I\);
the indicated minus sign on the weight-three ordinate makes the
lift have order three. Omitting that sign would instead give
the hyperelliptic involution as its cube. The actual lift
preserves \(v\), since \(y,A_0\) acquire the same multiplier.
It extends over all projective points, including \(s=-1,\infty\).

Rearranging \(v^2A_0+B_0=0\) gives
\begin{equation}
 v^2s^3+3s^2+(3-3v^2)s-v^2=0,\qquad
 \operatorname{disc}_s=81(v^4-v^2+1)^2 .
 \label{z2x-cubic-equation}
\end{equation}
The equation bounds the function-field degree by three, while
the three distinct powers of \(\rho\) bound it below by three.
Thus the fixed field is exactly \(\mathbb Q(v)\).

A fixed source satisfies \(s^2+s+1=0\). There \(A_0=-3s\),
\(B_0=-3\), and the two ordinates are nonzero. Moreover
\((s+1)^3=-1\), so both ordinates are fixed by the signed lift.
There are exactly four fixed geometric points. Their values
satisfy \(v^2=-1/s\), hence \(v^4-v^2+1=0\), giving four
distinct values. Away from fixed points the cyclic group acts
freely, so there is no further ramification. The three points
over \(v=0\) are the rational branch points at \(s=0,-1,\infty\);
those over \(v=\infty\) are the three nonrational zeros of \(A_0\).
The explicit discriminant in~\eqref{z2x-cubic-equation} follows
from the cubic discriminant formula and agrees with this complete
projective calculation.
\end{proof}

\begin{lemma}[Rational two-torsion and a full three-adic residue exclusion]
\label{z2x-torsion-local}
The classes \([P_0-P_\infty]\) and \([P_{-1}-P_\infty]\)
generate a subgroup \((\mathbb Z/2\mathbb Z)^2\) of \(J(\mathbb Q)\),
where \(J=\operatorname{Jac}(H)\). The whole primitive source residue
\(S\equiv T\not\equiv0\pmod3\) is absent from \(H(\mathbb Q_3)\).
\end{lemma}
\begin{proof}
The divisors of \(s\) and \(s+1\) are respectively
\(2P_0-2P_\infty\) and \(2P_{-1}-2P_\infty\). Neither class
vanishes, since a degree-one principal difference would force
genus zero. Their sum cannot vanish either: the functions with
sole pole at \(P_\infty\) of order at most two form the span of
\(1,s\). Indeed the smooth affine ring is \(\mathbb Q[s,y]\);
its unique expressions \(a(s)+b(s)y\) have even and odd pole
orders, respectively, since \(s,y\) have orders two and five.
They cannot cancel. A linear function of \(s\) cannot vanish
at both zero and minus one. This proves independence.

For a projective \(\mathbb Q_3\)-source choose integral \(S,T\)
with at least one a unit. If \(S=T+3K\) and \(T\) is a unit,
then
\[
 A_0=-3(T^3-9K^2T-9K^3),\qquad
 B_0/3=ST(S+T)\equiv2T^3\pmod3.
\]
The valuation of \(-A_0B_0\) is exactly two, but its quotient
by nine is \(2T^6\equiv-1\pmod3\), a nonsquare. The weight-three
scaling of \(Y\) proves that the test covers every deeper source
in that projective residue, including denominators.

In the other primitive residues \(A_0\) is a unit. Unless one
of \(S,T,S+T\) is zero, exactly one has positive three-adic
valuation, which must be odd for \(-A_0B_0\) to be a square.
The zero cases are the retained branch points. These are only
necessary local conditions. In particular \(S\equiv T\) is
already incompatible with the actual unramified root condition
\(3\nmid S^2+ST+T^2\); its exclusion alone supplies no global
rank or positive-locus result.
\end{proof}

\begin{theorem}[Finite reductions and the rational Jacobian]
\label{z2x-finite-jacobian}
The Jacobian \(J\) is \(\mathbb Q\)-simple, its rational torsion
is exactly \((\mathbb Z/2\mathbb Z)^2\), and its rank is even.
The complete reduction data are
\begin{equation}
\begin{array}{c|rr|l|r}
 p&\#H(\mathbb F_p)&\#H(\mathbb F_{p^2})&
 P_p(T)&P_p(1)\\ \hline
 5&3&25&T^4-3T^3+4T^2-15T+25&12\\
 7&9&37&T^4+T^3-6T^2+7T+49&52
\end{array}
\label{z2x-finite-table}
\end{equation}
\end{theorem}
\begin{proof}
Reduction at five and seven is good: the cubic discriminant
\(81\), its values \(-1,1\) at the two linear roots, their
difference, and the leading coefficient \(-3\) are units.
The odd degree also retains one smooth point at infinity.

For completeness, the exact finite certificate uses
\(\mathbb F_{25}=\mathbb F_5[u]/(u^2-2)\) and
\(\mathbb F_{49}=\mathbb F_7[u]/(u^2-3)\). Each relation is
irreducible. For every pair \(a+bu\), it evaluates the quintic
and counts all ordinates by the complete square-fiber table.
It separately enumerates all prime-field ordinates and adds
the unique point at infinity. All 25 and 49 abscissa fibers
are retained in the canonical result cited below. These
exhaustive finite counts give~\eqref{z2x-finite-table} by
\[
 a_1=p+1-N_p,\qquad
 a_2=\frac{N_{p^2}-p^2-1+a_1^2}{2},\qquad
 P_p(T)=T^4-a_1T^3+a_2T^2-pa_1T+p^2.
\]

If \(J\) had a proper nonzero abelian subvariety over \(\mathbb Q\),
it would be \(\mathbb Q\)-isogenous to two elliptic curves.
As in Section~\ref{qs-section}, good reduction at five passes
to the factors, and their Frobenius polynomials would give
\[
 P_5(T)=(T^2-aT+5)(T^2-bT+5),\qquad a,b\in\mathbb Z.
\]
The coefficients force \(a+b=3\), \(ab=-6\), so
\((a-b)^2=33\), impossible. This proves simplicity over
\(\mathbb Q\), without an absolute-simplicity assertion.

Prime-to-\(p\) torsion reduces injectively at a good prime.
The five-primary part injects into the order-52 group at seven
and vanishes. The seven-primary part injects into the order-12
group at five and vanishes. All other primary parts are bounded
by both orders, so the full torsion order divides four. The
four points in Lemma~\ref{z2x-torsion-local} attain this bound.

The norm-pullback identity for the cyclic quotient in
Proposition~\ref{z2x-cubic-map}, including ramified divisor
multiplicities, gives \(1+\rho_*+\rho_*^2=0\) on \(J\), since
the quotient has zero Jacobian. Thus \(\mathbb Q(\sqrt{-3})\)
acts on \(J(\mathbb Q)\otimes\mathbb Q\), making its rational
dimension even. This uses Mordell--Weil finite generation and
does not infer a rank from the reduction counts.
\end{proof}

\begin{lemma}[The point-set consequence of rank zero]
\label{z2x-rank-zero-points}
If \(\operatorname{rank}J(\mathbb Q)=0\), then
\[
 H(\mathbb Q)=\{P_0,P_{-1},P_\infty\}.
\]
\end{lemma}
\begin{proof}
The Abel--Jacobi map \(R\mapsto[R-P_\infty]\) is injective,
because a principal difference of distinct points would have
degree one. Under the hypothesis its image lies in the
four-element torsion group. Three elements are the images of
the three stated points. If the fourth were the image of \(R\),
then
\[
 R+P_\infty\sim P_0+P_{-1}.
\]
The divisor on the right is not equivalent to \(2P_\infty\),
by the pole-space argument in Lemma~\ref{z2x-torsion-local}.
The canonical class is \(2P_\infty\), so Riemann--Roch gives
\(h^0(P_0+P_{-1})=1\). Its only effective representative is
itself, which cannot equal \(R+P_\infty\). This excludes the
fourth element. The hypothesis will be discharged by the
actual descent in Section~\ref{zd2x-section}.
\end{proof}

\paragraph{Exact evidence and scope.}
The standard-library program
\texttt{next\_\allowbreak replay\_\allowbreak zeta\_\allowbreak
squared\_\allowbreak reductions.py} exhaustively certifies the
finite inputs. Its canonical JSON has SHA256
\texttt{66ce937b9a33a530\allowbreak0f08b128f630b3b47\allowbreak
ddaeee80b28b7f5a7\allowbreak7427ed30c06b6c}.
The author and both independent reviewers actually replayed
the certificate. No rank oracle or rational height cutoff is
used. The statements above are ordinary proofs with exact
finite arithmetic support, not full Lean formalizations.
Even a complete fixed-curve classification does not settle
ramified first norms, arbitrary residuals, varying indices,
or uniform point-height estimates for the parent problem.
